Write $B \cong \Mat_m(D)$, $D \in \{\R, \C, \HH\}$, $d = \dim_\R D$, with unique irreducible module $W = D^m$, $\dim_\R W = md$. $B$ acts faithfully on $W$, hence so does $A$; a faithful module of $A$, split as in the statement by \claimref{transport.classification-real-clifford-algebras}, contains both irreducibles $V_\pm$ (of $\dim_\R = kf$ each), since otherwise one of the central projectors $\varepsilon_\pm \neq 0$ would act as zero. So $W \cong V_+^{\,a} \oplus V_-^{\,b}$ as an $A$-module with $a, b \geq 1$ and
\[
md \;=\; (a + b)\,k f \;\geq\; 2kf .
\]
Moreover the right $D$-action on $W$ commutes with $A$, so $D$ embeds unitally in the commutant (strictly, the right action embeds $D^{\mathrm{op}}$; harmless, as $\R$, $\C$ and $\HH$ are each isomorphic to their opposites, and the dimension bounds below are opposite-invariant) $\mathrm{End}_A(W) \cong \Mat_a(\mathbb{F}) \oplus \Mat_b(\mathbb{F})$ (using $\mathrm{End}_{\Mat_k(\mathbb{F})}(\mathbb{F}^k) \cong \mathbb{F}$ as real algebras), hence, projecting to each summand, unitally in both $\Mat_a(\mathbb{F})$ and $\Mat_b(\mathbb{F})$.

Case $\mathbb{F} = \R$ ($f = 1$, $2k^2 = 2^n$). If $D = \R$: $m \geq 2k$ and $\dim B = m^2 \geq 4k^2 = 2^{n+1}$, with equality iff $m = 2k$, $B \cong \Mat_{2k}(\R)$. If $D = \C$: $\C \hookrightarrow \Mat_a(\R)$ forces $a \geq 2$, likewise $b \geq 2$, so $2m = (a+b)k \geq 4k$, $m \geq 2k$, and $\dim B = 2m^2 \geq 8k^2 = 2^{n+2}$. If $D = \HH$: $a, b \geq 4$, so $4m \geq 8k$, $\dim B = 4m^2 \geq 16k^2$. Case $\mathbb{F} = \HH$ ($f = 4$, $8k^2 = 2^n$). If $D = \HH$: $4m \geq 8k$ gives $m \geq 2k$ and $\dim B = 4m^2 \geq 16k^2 = 2^{n+1}$, equality iff $B \cong \Mat_{2k}(\HH)$. If $D = \C$: $\C \hookrightarrow \Mat_a(\HH)$ allows $a \geq 1$, but then $2m = 4(a+b)k \geq 8k$ gives $m \geq 4k$ and $\dim B = 2m^2 \geq 32k^2$. If $D = \R$: $m = 4(a+b)k \geq 8k$ and $\dim B = m^2 \geq 64k^2$. In every case the bound holds, and the equality case is exactly $\Mat_{2k}(\mathbb{F})$. Finally $\Mat_{2k}(\mathbb{F}) \cong \Mat_2(\Mat_k(\mathbb{F})) \cong \Mat_2(\Cl^0(p,q))$ by Lemma \claimref{transport.corner-even-subalgebra} and Theorem \claimref{transport.rule-2-ideal-fusion-sharp-form}.
